\section*{الفصل \ref{ch:b1:organic-redox} --- الأكسدة والاختزال في الكيمياء العضوية}

\begin{solution}{exo:b1:organic-redox:1}
الإيثانول: \ce{CH3} $-3$، \ce{CH2OH} $-1$. الإيثانال: \ce{CH3} $-3$، CHO $+1$.
حمض الإيثانويك: \ce{CH3} $-3$، COOH $+3$. البروبانون: \ce{CH3} $-3$ (مرتين)،
C=O $+2$. حمض الميثانويك: $+2$.
\end{solution}

\begin{solution}{exo:b1:organic-redox:2}
البروبان-1-أول: البروبانال، ثم حمض البروبانويك. البروبان-2-أول: البروبانون.
2-ميثيل البروبان-2-أول: لا تفاعل (يبقى اللون البرتقالي).
\end{solution}

\begin{solution}{exo:b1:organic-redox:3}
\ce{CH3COCH3 + 2H+ + 2e- -> CH3CH(OH)CH3};
\ce{CH3CHO + 2H+ + 2e- -> CH3CH2OH}.
\end{solution}

\begin{solution}{exo:b1:organic-redox:4}
\ce{NaBH4}: بوتان-1-أول، بوتان-2-أول، لا تفاعل، لا تفاعل (فالحمض
لا يُنزع منه إلا البروتون). \ce{LiAlH4}: بوتان-1-أول، بوتان-2-أول، بوتان-1-أول
وإيثانول، بوتان-1-أول.
\end{solution}

\begin{solution}{exo:b1:organic-redox:5}
\ce{5CH3CH(OH)CH3 + 2MnO4- + 6H+ -> 5CH3COCH3 + 2Mn^2+ + 8H2O}.
$n = 1.20/60.0 = \qty{0.0200}{mol}$؛ والبرمنغنات $\frac25 \times 0.0200 =
\qty{8.0e-3}{mol}$، أي \qty{400}{mL} من المحلول ذي التركيز \qty{0.020}{mol/L}.
\end{solution}

\begin{solution}{exo:b1:organic-redox:6}
سخّن الكحول مع \omterm{def:b1:nernst:couple}{المؤكسد} في دورق مجهَّز للتقطير، عند
درجة حرارة بين درجتي الغليان: فيتقطر البروبانال (\qty{48}{\celsius})
كلما تكوّن ويفلت من مزيد من الأكسدة؛ ويبقى البروبان-1-أول
(\qty{97}{\celsius}) في الدورق.
\end{solution}

\begin{solution}{exo:b1:organic-redox:7}
يهاجم زوج B--H من \ce{BH4-} كربون الكربونيل، ويذهب زوج $\pi$ للرابطة C=O
إلى الأكسجين؛ ويأخذ \omterm{def:b1:alcohol-activation:alkoxide}{الألكوكسيد} بروتونًا من الإيثانول. والهيدروجين
الموجود على كربون الناتج (\ce{CH3CH2OH}، أحد الاثنين على C1) يأتي
من الكاشف؛ والهيدروجين الموجود على الأكسجين يأتي من المذيب.
\end{solution}

\begin{solution}{exo:b1:organic-redox:8}
البوتان-2-أول \omterm{def:b1:stereochemistry-in-depth:stereocentre}{كيرالي}، لكن الكيتون المستوي يُهاجَم بالتساوي على
وجهيه: \omterm{def:b1:stereochemistry-in-depth:enantiomers}{خليط راسيمي}، غير نشط ضوئيًا.
\end{solution}

\begin{solution}{exo:b1:organic-redox:9}
\ce{3CH3CH2OH + 2Cr2O7^2- + 16H+ -> 3CH3COOH + 4Cr^3+ + 11H2O}: يتحول ثنائي الكرومات
البرتقالي (الكروم \babelsublr{$+$VI}) إلى الكروم(III) الأخضر.
\end{solution}

\begin{solution}{exo:b1:organic-redox:10}
2-ميثيل بوتان-2-أول: الإيثانال + \ce{CH3CH2MgBr} يعطي بوتان-2-أول (إضافة،
لا أكسدة ولا اختزال)؛ ثم أكسدة إلى البوتانون؛ ثم \ce{CH3MgBr} على البوتانون،
ثم المعالجة. أكسدة واحدة. البنتان-3-أول: البروبانال + \ce{CH3CH2MgBr} في
خطوة واحدة، لا أكسدة ولا اختزال.
\end{solution}

\begin{solution}{exo:b1:organic-redox:11}
1. احمِ الكيتون على هيئة \omterm{def:b1:acetals-protection:hemiacetal}{أسيتال} حلقي (الإيثان-1,2-ثنائي أول، حمض،
دين--ستارك). 2. يختزل \ce{LiAlH4} في إيثر جاف الإستر إلى
الكحول الأولي (والميثانول). 3. يزيل الحمض المائي \omterm{def:b1:acetals-protection:hemiacetal}{الأسيتال}:
5-هيدروكسي بنتان-2-ون، \ce{CH3CO(CH2)3OH}.
\end{solution}

\begin{solution}{exo:b1:organic-redox:12}
\ce{RCH2OH}: الكربون الحامل لمجموعة OH عند $-1$؛ \ce{RCOOH}: $+3$؛ أربعة إلكترونات، كما في
\ce{RCH2OH + H2O -> RCOOH + 4H+ + 4e-}. والميثان ($-4$) إلى ثنائي أكسيد الكربون
($+4$): ثمانية إلكترونات، \ce{CH4 + 2H2O -> CO2 + 8H+ + 8e-}.
\end{solution}

\begin{solution}{pb:b1:organic-redox:1}
\textbf{1.} C1 (الحامل لمجموعة OH): 0؛ ومجموعات \ce{CH2} الخمس: $-2$.
\textbf{2.} C1 (C=O): $+2$؛ ومجموعات \ce{CH2} الخمس: $-2$.
\textbf{3.} كربونا COOH: $+3$؛ ومجموعات \ce{CH2} الأربع: $-2$.
\textbf{4.} \babelsublr{C1 ($0 \to +3$)} وجاره \babelsublr{C6 ($-2 \to +3$)}، الذي
يصير الكربون الكربوكسيلي الثاني؛ وتنكسر الرابطة C1--C6.
\textbf{5.} $3 + 5 = 8$ إلكترونات.
\textbf{6.} حمض النتريك قوي بما يكفي لكسر الرابطة C--C المجاورة
للكربونيل، وهو ما لا تفعله المؤكسدات الأخف.
\textbf{7.} \ce{C6H12O + 3H2O -> C6H10O4 + 8H+ + 8e-}.
\textbf{8.} \babelsublr{$+$V} و\babelsublr{$+$I}.
\textbf{9.} \ce{2HNO3 + 8H+ + 8e- -> N2O + 5H2O}.
\textbf{10.} بالجمع، تُختصر ثمانية إلكترونات، و\ce{8H+} وثلاثة جزيئات ماء:
\ce{C6H12O + 2HNO3 -> C6H10O4 + N2O + 2H2O}.
\textbf{11.} 8.
\textbf{12.} $10^6/146.0 = \qty{6.85e3}{mol}$ من الحمض؛ $2 \times 6.85 \times
10^3 \times 63.0 = \qty{8.6e5}{g}$، أي \qty{0.86}{t} من حمض النتريك.
\textbf{13.} \ce{H2O2 + 2H+ + 2e- -> 2H2O}.
\textbf{14.} \ce{C6H12O + 4H2O2 -> C6H10O4 + 5H2O}.
\textbf{15.} $4 \times 6.85 \times 10^3 \times 34.0 = \qty{9.3e5}{g}$،
أي \qty{0.93}{t}.
\textbf{16.} الماء: لا أكسيد للنيتروجين.
\textbf{17.} ثابت التوازن الكبير لا يقول شيئًا عن السرعة؛
فالمركب \ce{H2O2} يتفاعل ببطء دون \omterm{def:b1:elementary-steps:catalyst}{حفّاز}، كما يتفكك هو نفسه ببطء.
\textbf{18.} يختزل الهكسانون الحلقي فيعيده إلى الهكسانول الحلقي:
\ce{4C6H10O + BH4- -> B(OC6H11)4-}، ثم
\ce{B(OC6H11)4- + 4H2O -> 4C6H11OH + B(OH)4-}.
\textbf{19.} \qty{6.85e3}{mol}.
\textbf{20.} \qty{6.85e3}{mol} من \ce{N2O} (جزيء لكل جزيء حمض).
\textbf{21.} $6.85 \times 10^3/0.95 \times 100.0 = \qty{7.2e5}{g}$، أي \qty{0.72}{t}.
\textbf{22.} $6.85 \times 10^3 \times 44.0 = \qty{3.0e5}{g}$: \textbf{\qty{0.30}{t}
من \ce{N2O} لكل طن من حمض الأديبيك}.
\end{solution}
